% Source: Weinberg GR, physical PDF pp. 430--433
%         (printed pp. 409--412).
% Coverage: Section 14.1, The Cosmological Principle; equations
%           (14.1.1)--(14.1.12).
% Figures/tables/footnotes: no figures, tables, footnotes, or reference markers.
% Status: source-reviewed and compile-clean; visually proofed.
% Uncertainties: the source's time-last component notation is converted to the
%                binding time-first convention below.
%
% MODERNIZATION CHECK: spacetime components are ordered time first; spatial
% indices i,j remain Latin and every Killing vector has vanishing 0-component.
% LITERAL-OPERATOR AUDIT: the plus sign joining the two sums in (14.1.8) and
% the minus sign in (14.1.10) were checked against physical PDF p. 432.

\subsection{The Cosmological Principle}\label{sec:14.1}

The Cosmological Principle is the hypothesis that the universe is spatially
homogeneous and isotropic. Before applying this principle, we shall have to
formulate our intuitive ideas of homogeneity and isotropy in precise
mathematical terms.

First, let us fix our attention on one particular spacetime coordinate system,
of the sort that might be used by terrestrial cosmographers. Spatial
coordinates \(x^i\) can be constructed with origin \(x^i=0\) at the center of
the Milky Way, with coordinate directions fixed by the line of sight from the
Milky Way to some typical distant galaxies, and with a scale of distances
defined by the apparent luminosities of distant galaxies, or other suitable
objects, as seen from the Milky Way. To define a time coordinate, it is
convenient to use the evolving universe itself as a clock. It is believed that
several cosmic scalar fields, such as the proper energy density \(\rho\), or
the black-body radiation temperature \(T_\gamma\) (see Chapter~15), are
everywhere decreasing monotonically; choose any one of these, say a scalar
\(S\), and let the time of any event be any definite decreasing function
\(t(S)\) of the chosen scalar, when and where the event occurs. (We shall have
to reopen the question of how to define the time when we consider a steady
state universe, in Section~\ref{sec:14.8}.) The coordinates
\(x^\mu=(t,\mathbf{x})\) so defined will be called the \emph{cosmic standard
coordinate system}.

The Cosmological Principle can be formulated as a statement about the
existence of equivalent coordinate systems. Suppose that we use the cosmic
standard coordinate system to carry out astronomical observations, determining
(never mind how!) the metric tensor \(g_{\mu\nu}\), the energy-momentum tensor
\(T_{\mu\nu}\), and all other cosmic fields, as functions of the cosmic standard
coordinates \(x^\mu\). A different set of spacetime coordinates \(x'^\mu\) may
be considered equivalent to the cosmic standard coordinates if the whole
history of the universe appears the same in the \(x'^\mu\) coordinate system
as in the cosmic standard coordinate system. This requires that every cosmic
field \(g'_{\mu\nu}(x')\), \(T'_{\mu\nu}(x')\), and so on, must be the same
function of the \(x'^\mu\) as the corresponding quantities
\(g_{\mu\nu}(x)\), \(T_{\mu\nu}(x)\), and so on, are of the standard
coordinates \(x^\mu\). That is, at any coordinate point \(y^\mu\), we must
have
\begin{equation}
  g'_{\mu\nu}(y)=g_{\mu\nu}(y),
  \tag{14.1.1}\label{eq:14.1.1}
\end{equation}
\begin{equation}
  T'_{\mu\nu}(y)=T_{\mu\nu}(y),\qquad\text{etc.}
  \tag{14.1.2}\label{eq:14.1.2}
\end{equation}
In the language of the last chapter, Eq.~\eqref{eq:14.1.1} says that the
coordinate transformation \(x\longrightarrow x'\) must be an isometry, and
Eq.~\eqref{eq:14.1.2} says that \(T_{\mu\nu}\), and so on, must be
form-invariant under this transformation.

In particular, Eq.~\eqref{eq:14.1.2} will have to hold for the scalar \(S\)
used to define our cosmic standard time \(t\). Since \(S\) is by definition a
function only of \(t\), and a scalar, Eq.~\eqref{eq:14.1.2} for \(S\) reads,
at \(y=x'\),
\[
  S(t')=S'(x')=S(x)=S(t),
\]
and so
\begin{equation}
  t'=t.
  \tag{14.1.3}\label{eq:14.1.3}
\end{equation}
All coordinate systems that are equivalent to the cosmic standard system
necessarily use cosmic standard time.

The assumption of spatial isotropy can now be formulated as the requirement
that there exists a family of coordinate systems \(x'^\mu(x;\theta)\),
depending on three independent parameters \(\theta^1,\theta^2,\theta^3\),
which are equivalent to the cosmic standard coordinates, and which have the
same origin, that is,
\begin{equation}
  x'^i(0,t;\theta)=0.
  \tag{14.1.4}\label{eq:14.1.4}
\end{equation}
We can intuitively think of the three parameters \(\theta^i\) as Euler angles
that specify the orientation of the \(x'^i\) coordinate axes relative to the
\(x^i\) coordinate axes, but it is unnecessary to be so specific; the
important thing is that there be three independent parameters. (In
formulating this assumption, we have tacitly assumed that the privileged
Lorentz frame in which the universe appears isotropic happens to coincide more
or less with our own galaxy.)

It is a little trickier to formulate the assumption of homogeneity. Clearly,
homogeneity does not mean that any object can be chosen as the origin of a
coordinate system equivalent to our cosmic standard coordinates---after all,
the universe looks different to an observer moving away from the Milky Way at
half the speed of light than it does to us! The most we can expect is that
every point \(x^\mu\) in spacetime is on some ``fundamental trajectory''
\(x^i=X^i(t)\), which can serve as the origin of a coordinate system
\(x'^\mu\) equivalent to the cosmic standard system. (This is closely related
to a postulate called Weyl's principle, used in some formulations of
cosmology.) The Milky Way appears to be a rather ordinary galaxy, more or less
at rest with respect to its nearest neighbors, so we can expect that the
fundamental trajectories \(X(t)\) are pretty well defined by the motions of
typical members of the cosmic gas of galaxies, but this is by no means an
essential part of the assumption of homogeneity. The important point is that,
since the \(X(t)\) at any time \(t\) fill up all space, they are determined by
three independent parameters \(a^i\), which can be taken, for instance, as the
values \(a^i=X^i(T)\) of \(X^i(t)\) at some particular time \(t=T\). Thus
homogeneity means that there is a three-parameter set of coordinates
\(\bar x'^\mu(x;a)\), which are equivalent to the cosmic standard coordinates
and which have origin on the trajectory \(x^i=X^i(t;a)\), that is,
\begin{equation}
  \bar x'^i\bigl(X(t;a),t;a\bigr)=0.
  \tag{14.1.5}\label{eq:14.1.5}
\end{equation}
To be more precise, the \(X(t;a)\) are the trajectories of the privileged
observers to whom the universe appears isotropic.

Putting this together, we see that the Cosmological Principle entails the
existence of two independent three-parameter families of coordinate
transformations \(x\longrightarrow x'\), \(x\longrightarrow\bar x'\), which
are isometries in the sense of Eq.~\eqref{eq:14.1.1}, and which according to
Eq.~\eqref{eq:14.1.3} leave the time coordinate invariant. The universe
therefore satisfies the requirements assumed in Section~\ref{sec:13.5} for a
four-dimensional space with three-dimensional maximally symmetric subspaces
\(t=\text{const}\).

(To see this in detail, we can descend to the case of infinitesimal
transformations, letting \(\theta^i\) and \(a^i\) approach zero. There are then
six ``Killing vectors'' \(\zeta_j^\mu(x)\) and \(\xi_j^\mu(x)\), defined by
\begin{equation}
  \zeta_j^i(x)
  \equiv
  \left.
  \frac{\partial x'^i(x;\theta)}{\partial\theta^j}
  \right|_{\theta=0},
  \qquad
  \zeta_j^0(x)\equiv0,
  \tag{14.1.6}\label{eq:14.1.6}
\end{equation}
\begin{equation}
  \xi_j^i(x)
  \equiv
  \left.
  \frac{\partial\bar x'^i(x;a)}{\partial a^j}
  \right|_{a=0},
  \qquad
  \xi_j^0(x)\equiv0.
  \tag{14.1.7}\label{eq:14.1.7}
\end{equation}
It is only necessary to show that these six vectors are independent. Suppose
that they satisfy a linear relation
\begin{equation}
  \sum_j c^j(t)\zeta_j^i(x)
  +\sum_j\bar c^{\,j}(t)\xi_j^i(x)=0.
  \tag{14.1.8}\label{eq:14.1.8}
\end{equation}
At the origin, Eqs.~\eqref{eq:14.1.4} and \eqref{eq:14.1.5} give
\begin{equation}
  \zeta_j^i(0,t)=0,
  \tag{14.1.9}\label{eq:14.1.9}
\end{equation}
\begin{equation}
  \xi_j^i(0,t)
  =
  -\left.
  \frac{\partial X^i(t;a)}{\partial a^j}
  \right|_{a=0},
  \tag{14.1.10}\label{eq:14.1.10}
\end{equation}
so at \(x^i=0\), Eq.~\eqref{eq:14.1.8} gives
\[
  \left.
  \sum_j\bar c^{\,j}(t)\frac{\partial X^i(t;a)}{\partial a^j}
  \right|_{a=0}=0.
\]
Since the \(a^i\) are independent parameters, this requires that
\begin{equation}
  \bar c^{\,j}(t)=0.
  \tag{14.1.11}\label{eq:14.1.11}
\end{equation}
Going back to Eqs.~\eqref{eq:14.1.8} and \eqref{eq:14.1.6}, we then have
\[
  \left.
  \sum_j c^j(t)
  \frac{\partial x'^i(x;\theta)}{\partial\theta^j}
  \right|_{\theta=0}=0,
\]
and since the \(\theta^i\) are independent parameters, this requires that
\begin{equation}
  c^j(t)=0.
  \tag{14.1.12}\label{eq:14.1.12}
\end{equation}
Thus there are six independent Killing vectors with vanishing time component,
the maximum number possible (see Section~\ref{sec:13.1}) in three
dimensions.)

In summary, the Cosmological Principle can be formulated in the language of
Chapter~13 as follows:

\begin{enumerate}
\item[(i)] The hypersurfaces with constant cosmic standard time are maximally
  symmetric subspaces of the whole of spacetime.

\item[(ii)] Not only the metric \(g_{\mu\nu}\), but all cosmic tensors such as
  \(T_{\mu\nu}\), are form-invariant with respect to the isometries of these
  subspaces.
\end{enumerate}
